\section{Agent interface, implementation, and non-functional properties}
\label{sec:implementation}
\label{sec:impl}

\subsection{The tool surface}
\label{sec:impl-tools}

The platform is exposed to agents as a single Model Context Protocol
server~\cite{mcp} carrying 95 tools in twelve groups: knowledge-base catalog;
document read/write; file system; parsing and task status; tags; search; the
verification gate; vector index and repair; knowledge graph; experience;
persona; and platform lifecycle and health. The count is the number of registered tool decorators in
\texttt{kb-mcp/server.py}, checked one-to-one against the function
definitions they register; the project's own documentation had drifted to a
lower number, so we report the measured one. Of the twelve groups, the paper
evaluates seven (catalog, documents, parsing, search, gate, index and graph);
the persona subsystem and the platform-lifecycle tools are outside its scope,
and we say so rather than claiming coverage of the whole surface. The interface is not a reduced
agent-only view: it owns no private data path. Writes go through the same
HTTP services the user interface calls, reads read the same
\texttt{.tree-fs.json} and \texttt{.knowledge-base.yml} registries the web
server reads, and the MCP process itself contains no HTTP code, only a
client module.

Three governance rules make a surface this wide safe to hand to a model.
First, read-path tools and write-path tools are disjoint families, and
retrieval lanes are granted the read family only. Second, a stored parsed
document cannot enter a knowledge base through the generic document writer:
it must pass the parsed-document writer (\texttt{kb\_doc\_save\_parsed}),
the single entry point that carries parse provenance and forces the
ingestion quality gates to run. Third, long-running operations---a parse
batch, a re-index, a graph build, a meditation pass (the scheduled pass that
rewrites the experience store's rolling summary from accumulated questions,
retrievals and feedback)---return a task id
immediately instead of blocking the session, and are polled through the
task-status tools backed by an in-process task registry.

\subsection{Harness integration and lane policy}
\label{sec:impl-harness}

One chat endpoint adapts fourteen agent harnesses to a uniform message
stream, so a vendor CLI, an ACP-compatible agent, and the platform's own
engine all reach the same retrieval lanes. The lane policy of
Section~\ref{sec:lanes} is enforced here rather than by prompt discipline:
each lane is attached to a tool allow-list, of which 21 entries are
read-only MCP retrieval tools, and to a deny-list of more than 40
harness-internal tools (file editors, web search, sub-agent spawning). The
gate was observed working in the harness campaign: on its first turn a model
attempted to end the session early, the attempt was blocked by the
deny-list, and the retrieval that followed was unaffected. Sessions are
budgeted as well, with per-lane turn and call caps
(Table~\ref{tab:lanes}); the campaign ran 18 end-to-end scenarios, all of
which passed with no fabricated answer.

Three operational findings from that campaign matter to callers. A trivial
ping prompt triggered a six-skill pre-flight before any retrieval, and the
probe was killed by a 30\,s timeout; the platform's default caller timeout
is therefore 420\,s, not the few seconds a request such as ``list my
libraries'' suggests. The streaming endpoint emits bare \texttt{data:}
frames with no event names, so a client must parse data-only frames instead
of waiting for a named event. Finally, engine choice dominates cost by one
to two orders of magnitude at the same quality tier: \$0.90--1.80 per
question through the agentic surface, \$0.11--0.36 through the Claude
harness, and \$0.03--0.07 through the OMP harness (the three rows of
Table~\ref{tab:e5-cost}). Wall clock was 50--176\,s, 23--52\,s, and
34--171\,s respectively, but the third engine's engine-side time was only
15--83\,s: its extra wall clock is harness cold start of 50--90\,s per call,
not retrieval. One residual gap appeared in the same campaign: a document
ingested through the third harness had no graph node, while the same
operation through the agent surface created one.

\subsection{Implementation and runtime engineering}
\label{sec:impl-runtime}

The backend is Python 3.12 with FastAPI. ChromaDB~\cite{chroma} stores one
HNSW~\cite{hnsw} collection per knowledge base in cosine space; embeddings
are BGE-M3~\cite{bge-m3}, 1024-dimensional and normalized, encoded in
batches of 32 by a lazy-loaded singleton with an LRU encode cache. Neo4j
5.20~\cite{neo4j} holds the document graph, jieba~\cite{jieba} tokenizes the
keyword lane, MinerU~\cite{mineru} performs parsing and OCR, and MCP is
spoken over stdio and SSE.

The enforcement described in Section~\ref{sec:lanes} depends on engineering
invisible in the interfaces. The judging model runs as a resident worker
behind a machine-level file lock, with an LRU verdict cache and explicit
CUDA-cache release between workloads. A cold judge load costs 10--12\,s per
batch; the resident worker answers in 1--2\,s. The worker also removes a
failure mode met in production: when judge and bulk-embedding workloads ran
interleaved without cache cleanup, one end-to-end run degraded from 68\,s to
27 minutes under the Windows display driver model's paging behaviour. Judge
calls are serialized by one lock, because two concurrent judgments on a
single GPU raise a native allocator failure; a call is held to a 900\,s
ceiling and fails closed beyond it. MinerU
is launched as a hidden subprocess with platform-specific parent-death
cleanup---a Windows Job Object with kill-on-close, Linux
\texttt{prctl(PR\_SET\_PDEATHSIG)}, and a process group plus
\texttt{atexit} on macOS---and its stdout is redirected to a log file rather
than a pipe, because the pipe variant failed with a spurious
\texttt{Errno 22} once the reader closed. Before start, the backend probes
its own port with a bare socket bind and refuses to start if the port is
occupied, so a surviving process cannot be shadowed by a second one. Vector
writes are serialized per collection, after one production incident traced
to concurrent segment handles; the keyword index updates incrementally.
Every HTTP client is constructed with \texttt{trust\_env=False}, so a
corporate proxy cannot intercept calls to localhost.

\subsection{Configuration, deployment, and reproducibility}
\label{sec:impl-deploy}

A single \texttt{config.yml} is the source of truth for ports and
thresholds; environment variables override it. The mechanism is used in
practice, not only documented: in the harness campaign the services ran on
web 6789, backend 8771, and kb-mcp 8000, the backend port being an override
of the shipped development default. Rate limiting (a general bucket and a
smaller bucket for expensive routes), authentication tokens for the MCP,
user, and legacy clients, and the retrieval thresholds all live in the same
file. The \texttt{ragctl} command-line tool covers setup, start, status,
dependency, configuration, and update operations. The system runs
self-hosted on a single workstation: ChromaDB runs in-process, Neo4j runs
locally on bolt port 7687, and MinerU runs as a local subprocess, so no
external service is required for a deployment. The same file carries the
higher retrieval threshold (0.55) used by the experience subsystem, whose
rolling summary of operational lessons is regenerated by a rules-based pass
every fifth ingestion. That loop is implemented but not evaluated:
Section~\ref{sec:eval}'s experience evidence is one scenario, and scheduled
distillation at scale, approval, decay and application feedback remain
unexercised (Section~\ref{sec:e3}). Source code and all experiment scripts are public at
\url{https://github.com/kingdol666/rag-knowledge}.
